\documentclass[10pt,a4paper]{amsart}
\usepackage{amsmath,amssymb,amsthm,graphicx}
\usepackage{verbatim}
\newenvironment{reference}{\comment}{\endcomment}
%\newenvironment{reference}{}{}
\newenvironment{slogan}{\comment}{\endcomment}
\newenvironment{history}{\comment}{\endcomment}

% For commutative diagrams we use Xy-pic
\usepackage[all]{xy}

% We use 2cell for 2-commutative diagrams.
\xyoption{2cell}
% Theorem environments.
%
\theoremstyle{plain}
\newtheorem{theorem}[subsection]{Theorem}
\newtheorem{proposition}[subsection]{Proposition}
\newtheorem{lemma}[subsection]{Lemma}

\theoremstyle{definition}
\newtheorem{definition}[subsection]{Definition}
\newtheorem{example}[subsection]{Example}
\newtheorem{exercise}[subsection]{Exercise}
\newtheorem{situation}[subsection]{Situation}

\theoremstyle{remark}
\newtheorem{remark}[subsection]{Remark}
\newtheorem{remarks}[subsection]{Remarks}

\numberwithin{equation}{subsection}

% Macros
%
\def\lim{\mathop{\mathrm{lim}}\nolimits}
\def\colim{\mathop{\mathrm{colim}}\nolimits}
\def\Spec{\mathop{\mathrm{Spec}}}
\def\Hom{\mathop{\mathrm{Hom}}\nolimits}
\def\Ext{\mathop{\mathrm{Ext}}\nolimits}
\def\SheafHom{\mathop{\mathcal{H}\!\mathit{om}}\nolimits}
\def\SheafExt{\mathop{\mathcal{E}\!\mathit{xt}}\nolimits}
\def\Sch{\mathit{Sch}}
\def\Mor{\mathop{\mathrm{Mor}}\nolimits}
\def\Ob{\mathop{\mathrm{Ob}}\nolimits}
\def\Sh{\mathop{\mathit{Sh}}\nolimits}
\def\NL{\mathop{N\!L}\nolimits}
\def\CH{\mathop{\mathrm{CH}}\nolimits}
\def\proetale{{pro\text{-}\acute{e}tale}}
\def\etale{{\acute{e}tale}}
\def\QCoh{\mathit{QCoh}}
\def\Ker{\mathop{\mathrm{Ker}}}
\def\Im{\mathop{\mathrm{Im}}}
\def\Coker{\mathop{\mathrm{Coker}}}
\def\Coim{\mathop{\mathrm{Coim}}}

% Boxtimes
%
\DeclareMathSymbol{\boxtimes}{\mathbin}{AMSa}{"02}

%
% Macros for moduli stacks/spaces
%
\def\QCohstack{\mathcal{QC}\!\mathit{oh}}
\def\Cohstack{\mathcal{C}\!\mathit{oh}}
\def\Spacesstack{\mathcal{S}\!\mathit{paces}}
\def\Quotfunctor{\mathrm{Quot}}
\def\Hilbfunctor{\mathrm{Hilb}}
\def\Curvesstack{\mathcal{C}\!\mathit{urves}}
\def\Polarizedstack{\mathcal{P}\!\mathit{olarized}}
\def\Complexesstack{\mathcal{C}\!\mathit{omplexes}}
% \Pic is the operator that assigns to X its picard group, usage \Pic(X)
% \Picardstack_{X/B} denotes the Picard stack of X over B
% \Picardfunctor_{X/B} denotes the Picard functor of X over B
\def\Pic{\mathop{\mathrm{Pic}}\nolimits}
\def\Picardstack{\mathcal{P}\!\mathit{ic}}
\def\Picardfunctor{\mathrm{Pic}}
\def\Deformationcategory{\mathcal{D}\!\mathit{ef}}
\usepackage{cleveref}
\crefname{originalreference}{original reference}{original references}
\Crefname{originalreference}{Original reference}{Original references}
\makeatletter
\newlabel{example-dgm-dg-cat}{{09LC}{source}{}{sourcepointer-1}{}}
\newlabel{example-dgm-dg-cat@cref}{{[originalreference][0][]09LC}{[1][0][]source}{}{}{}}
\makeatother
\crefname{originalreference}{original reference}{original references}
\Crefname{originalreference}{Original reference}{Original references}
\makeatletter
\newlabel{lemma-property-P-sequence}{{09KL}{source}{}{sourcepointer-2}{}}
\newlabel{lemma-property-P-sequence@cref}{{[originalreference][0][]09KL}{[1][0][]source}{}{}{}}
\makeatother
\title{Derived Hom against a compact bimodule with representable dual is derived tensor with the dual}
\author{The Stacks Project Authors}
\date{stacks-project:dga.tex:local-snapshot-20261012; GNU FDL 1.2 or later}
\begin{document}
\maketitle
\noindent\textit{Editorial note.} All delivered mathematical text is the author's own, copied from the pinned Stacks Project chapter file sources/stacks-dga/pinned-original-source.tex; no statement or proof step was rewritten, reordered or removed.\par\medskip
\noindent\textit{Editorial note.} The counted proof is the author's complete argument: the (B, A)-bimodule structure on S from the dg-category composition, the comparison map c\_N and its functoriality, the derived transformation through property-(P) and property-(I) replacements, and the two-step compactness reduction to N = B[k] closed by assumption (2) are all retained verbatim.\par\medskip
\noindent\textit{Editorial note.} The label lemma-RHom-is-tensor names the counted statement, and the proof environment immediately following it in the source is delivered verbatim as its proof.\par\medskip
\noindent\textit{Editorial note.} The counted statement is carried by the author's own lemma environment in the chapter source.\par\medskip
\noindent\textit{Allowed standard prerequisites (editorial).}\begin{itemize}
\item The differential graded category structure and bimodule actions. Same chapter, Example example-dgm-dg-cat and the module structures of section-restriction. Supplies the composition on Hom complexes that gives S = Hom(T, B) its (B, A)-bimodule structure and the counit-type multiplication defining the comparison map c\_N.
\item Compactness and the two-term filtrations of property (P). Same chapter, Lemma lemma-property-P-sequence; the compact-object formalism of proposition-compact. Lets both sides, being exact functors of triangulated categories with T compact, be checked on finite filtrations and finally on shifts of B.
\end{itemize}\medskip
\ifdefined\hypertarget\hypertarget{sourcepointer-1}{}\fi
\noindent\textit{Source reference.} Original reference 09LC: sources/stacks-dga/pinned-original-source.tex, source lines 3525--3525. Prerequisite pointer invoked at proof line(s) 5806; it names a standard fact of the same chapter that the counted argument uses, and every step of the counted proof built on it is retained verbatim in the delivered body.\par\medskip
\ifdefined\hypertarget\hypertarget{sourcepointer-2}{}\fi
\noindent\textit{Source reference.} Original reference 09KL: sources/stacks-dga/pinned-original-source.tex, source lines 2425--2425. Prerequisite pointer invoked at proof line(s) 5827; it names a standard fact of the same chapter that the counted argument uses, and every step of the counted proof built on it is retained verbatim in the delivered body.\par\medskip
\begin{lemma}
\label{lemma-RHom-is-tensor}
Let $R$ be a ring.
Let $(A, \text{d})$ and $(B, \text{d})$ be differential graded $R$-algebras.
Let $T$ be a differential graded $(A, B)$-bimodule.
Assume
\begin{enumerate}
\item $T$ defines a compact object of $D(B, \text{d})$, and
\item $S = \Hom_{\text{Mod}^{dg}_{(B, \text{d})}}(T, B)$
represents $R\Hom(T, B)$ in $D(A, \text{d})$.
\end{enumerate}
Then $S$ has a structure of a differential graded $(B, A)$-bimodule
and there is an isomorphism
$$
N \otimes_B^\mathbf{L} S \longrightarrow R\Hom(T, N)
$$
functorial in $N$ in $D(B, \text{d})$.
\end{lemma}

\begin{proof}
Write $\mathcal{B} = \text{Mod}^{dg}_{(B, \text{d})}$.
The right $A$-module structure on $S$ comes from the map
$A \to \Hom_\mathcal{B}(T, T)$ and the composition
$\Hom_\mathcal{B}(T, B) \otimes \Hom_\mathcal{B}(T, T)
\to \Hom_\mathcal{B}(T, B)$ defined in Example \ref{example-dgm-dg-cat}.
Using this multiplication a second time there is a map
$$
c_N :
N \otimes_B S = \Hom_\mathcal{B}(B, N) \otimes_B \Hom_\mathcal{B}(T, B)
\longrightarrow
\Hom_\mathcal{B}(T, N)
$$
functorial in $N$. Given $N$ we can choose quasi-isomorphisms
$P \to N \to I$ where $P$, resp.\ $I$ is a differential graded $B$-module
with property (P), resp.\ (I). Then using $c_N$ we obtain a map
$P \otimes_B S \to \Hom_\mathcal{B}(T, I)$ between the objects
representing $S \otimes_B^\mathbf{L} N$ and $R\Hom(T, N)$.
Clearly this defines a transformation of functors $c$ as in the lemma.

\medskip\noindent
To prove that $c$ is an isomorphism of functors, we may
assume $N$ is a differential graded $B$-module which
has property (P). Since $T$ defines a compact object in
$D(B, \text{d})$ and since both sides of the arrow define
exact functors of triangulated categories, we reduce using
Lemma \ref{lemma-property-P-sequence}
to the case where $N$ has a finite filtration whose
graded pieces are direct sums of $B[k]$. Using again that
both sides of the arrow are exact functors of triangulated
categories and compactness of $T$ we reduce to
the case $N = B[k]$. Assumption (2) is exactly the
assumption that $c$ is an isomorphism in this case.
\end{proof}

\end{document}
